\section{Introduction}



Live event streaming services such as Prime Video operate under stringent reliability requirements, particularly during large-scale broadcasts like NFL Thursday Night Football (TNF) and Premier League matches, where millions of viewers depend on uninterrupted service. Even a single faulty code change can cause widespread disruptions that are difficult to mitigate during events. To reduce this risk, organizations commonly employ broad deployment freezes during high-traffic windows. While these freezes prevent failures, they delay feature rollouts, create pending change accumulations, and require special handling for urgent updates. As the number and complexity of live events continue to grow, the operational cost of such freezes increases correspondingly, making more precise assessment of individual code-change risk increasingly important.

The central challenge is determining which changes warrant caution without relying on coarse, event-long freezes. Although prior research on change- and defect-prediction provides useful insights~\cite{kamei2012large,6772130}, many approaches depend on developer-specific metadata or substantial historical project information. These dependencies raise privacy considerations, hinder adoption across diverse teams, and limit applicability to new or rapidly evolving codebases. Other models demonstrate effectiveness only within a single dataset, providing limited evidence of generalizability across languages, repositories, or organizational environments.

These limitations motivate an approach using only code modification information. Diff-level characteristics avoid sensitive metadata, apply immediately to new projects including cold-start scenarios such as new code bases, and behave consistently across heterogeneous environments. Such a model would be particularly valuable in contexts like Prime Video's live event operations, where accurate and lightweight risk assessment is needed to make informed deployment decisions during TNF and other high-visibility broadcasts.

In this work, we investigate the predictive utility of diff-aware features for assessing risky code changes. Our intended deployment model is as a \emph{guardrail}: the system flags potentially risky changes for human review rather than autonomously blocking deployments, keeping engineers in the loop for final deployment decisions. This human-in-the-loop design motivates our preference for recall over precision: it is preferable to surface a change for review unnecessarily than to miss a genuinely risky one, leading us to select a threshold that weights recall more (Section~\ref{subsec:classification}). By constructing a comprehensive set of quantitative and qualitative attributes derived directly from code modifications, and evaluating them across both industrial and public datasets, we assess whether code-change information alone can support practical and reliable risk prediction. Our results indicate that a well-curated set of diff-level features captures meaningful risk signals, enabling more targeted and efficient deployment decision-making in environments where reliability and development velocity must be jointly balanced.

Our contributions are: (1) a diff-aware risk assessment framework applicable to cold-start scenarios without developer metadata, (2) LLM-based multi-language feature extraction that replaces brittle language-specific AST parsers, and (3) empirical evidence that LLMs are effective for structured feature extraction but lag behind feature-engineered models for end-to-end risk classification.

The remainder of this paper is organized as follows: Section~2 reviews related work, Section~3 describes our methodology including data collection, feature extraction, and classification, Section~4 presents the classification and ablation results, Section~5 discusses industrial deployment considerations and feature insights, Section~6 addresses threats to validity, and Section~7 concludes with future directions.
